\documentclass[18pt]{article}
\usepackage{geometry}
\geometry{top=2cm}
\usepackage{amsmath,amsfonts,amssymb}
\usepackage{natbib}
\usepackage{booktabs}
\usepackage[utf8]{inputenc}
\usepackage[T1]{fontenc}
\usepackage{amssymb}
\usepackage[version=4]{mhchem}
\usepackage{stmaryrd}
\usepackage{bbold}
\usepackage{nopageno}
\usepackage{hyperref}
\usepackage{bookmark}
\usepackage{accents}
\author{Adam Handwerger}
\pagestyle{plain}
\begin{document}
\begin{flushleft}
    1 Kernels\\
    \vspace{.4cm}
    1. Let $\chi=\mathbb{R}^{+} $\\
    \vspace{.4cm}
    (a) Verify that $ \min(x,y)=\int_{0}^\infty (\mathbb{I}_{t<x}) (\mathbb{I}_{t<y}) \,dt $\\
    \vspace{.4cm}
    Choose $x,y \in \mathbb{R}^{+}$.\\
    \vspace{.4cm}
    Case 1, $x \geq y$\\
    \vspace{.4cm}
    $\int_{0}^{\infty} (\mathbb{I}_{t<x}) (\mathbb{I}_{t<y}) \,dt=\int_{0}^{y} (\mathbb{I}_{t<x}) (\mathbb{I}_{t<y}) \,dt+$\\
    $\int_{y}^{\infty} (\mathbb{I}_{t<x}) (\mathbb{I}_{t<y})=\int_{0}^{y} (1)(1) \,dt +0= \int_{0}^{y} \,dt=y$\\
    \vspace{.4cm}
    Case 2, $ x\leq y$\\
    \vspace{.4cm}
    $\int_{0}^{\infty} (\mathbb{I}_{t<x}) (\mathbb{I}_{t<y}) \,dt=\int_{0}^{x} (\mathbb{I}_{t<x}) (\mathbb{I}_{t<y}) \,dt+$\\
    $\int_{x}^{\infty} (\mathbb{I}_{t<x}) (\mathbb{I}_{t<y})=\int_{0}^{x} (1)(1) \,dt +0= \int_{0}^{x} \,dt=x$\\
    \vspace{.4cm}
    Combining both cases gives the desired result.\\
    \vspace{.4cm}
    (b) Show that $K(x,y)=\min(x,y)$ is pos. definite over $\mathbb{R}^{+}$.\\
    \vspace{.4cm}
    Case 1, $x_i \leq x_j$ and $\alpha\in\mathbb{R}^n$\\
    \vspace{.4cm}
    In the same manner as above\\
    \vspace{.4cm}
    $\sum_{ij}^n \alpha_i \alpha_j \int_{0}^{\infty} (\mathbb{I}_{t<x_i}) (\mathbb{I}_{t<x_j}) \,dt=\int_{0}^{x_i} (\sum_i \alpha_i \mathbb{I}_{t<x}) (\sum_j \alpha_j \mathbb{I}_{t<y}) \,dt=$\\
    $\int_0^{x_i} (\sum_i \alpha_i)(\sum_j \alpha_j) \,dt=(\sum_k \alpha_k)^2 x_i \geq 0$ for $x_i \in \mathbb{R}^{+}$.\\
    \vspace{.4cm}
    Likewise for $ x_i \geq x_j$,\\
    \vspace{.4cm}
    $\sum_{ij}^n \alpha_i \alpha_j \int_{0}^{\infty} (\mathbb{I}_{t<x_i}) (\mathbb{I}_{t<x_j}) \,dt=\int_{0}^{y_i} (\sum_i \alpha_i \mathbb{I}_{t<x}) (\sum_j \alpha_j \mathbb{I}_{t<y}) \,dt=$\\
    $\int_0^{y_i} (\sum_i \alpha_i)(\sum_j \alpha_j) \,dt=(\sum_k \alpha_k)^2 y_i \geq 0$ for $y_i \in \mathbb{R}^{+}$.\\
    \vspace{.4cm}
    Hence K is pos. definite.\\
    \vspace{.4cm}
    (c) Show that $\max(x,y)$ is not pos. definite over $\mathbb{R}^{+}$.
    \vspace{.4cm}
    Consider the case when n=2. Let $v=(\alpha_1,\alpha_2)$ and $\gamma=\max(x_1,x_2)$. Then\\
    \end{flushleft}
    $$ v^T K v=
    \begin{bmatrix} \alpha_1 & \alpha_2 \end{bmatrix}
        \begin{bmatrix} x_1 & \gamma\\
                        \gamma & x_2 \end{bmatrix}
    \begin{bmatrix} \alpha_1 \\ \alpha_2 \end{bmatrix}\\
    $$
    $\hspace{4.5cm} =\alpha_1^2 x_1 + 2 \alpha_1 \alpha_2 \gamma + \alpha_2^2 x_2$\\
    \vspace{.4cm}
    \\
    Since this must hold for any $v \in \mathbb{R}^2$, choose $\alpha_1=-\alpha_2$.\\
    \\
    Then the expression above becomes\\
    \begin{flushleft}
    $\alpha_1^2 x_1 + 2 \alpha_1 \alpha_2 \gamma + \alpha_2^2 x_2 = \alpha_1^2 (x_1 + x_2 -2\gamma)=$\\
    $\alpha_1^2(x_1+x_2-2 \max(x_1,x_2)) \leq 0$\\
\vspace{.4cm}
2. Let $(\Omega,\mathcal{A}, P)$ be a probability space\\
\vspace{.4cm}
(a) Show that for any two events $A,B\in \mathcal{A}$,\\
$K(A,B)=P(A \cap B)$ is pos. definite.\\
\vspace{.4cm}
Choose $(\alpha_1,\dots,\alpha_n) \in \mathbb{R}^n$. Then\\
\vspace{.4cm}
$\sum_{ij}^n \alpha_i \alpha_j K(x_i,x_j)=\sum_{ij}^n \alpha_i \alpha_j P(A_i \cap A_j)=$\\
$\sum_{ij}^n \alpha_i \alpha_j E(\mathbb{I}_{A_i} \mathbb{I}_{A_j})=E\left( \left(\sum_i \alpha_i \mathbb{I}_{A_i}\right)\left(\sum_j \alpha_j \mathbb{I}_{A_j}\right)\right)=$\\
$E\left( \left(\sum_k \alpha_k \mathbb{I}_{A_k} \right)^2 \right) \geq 0$\\
\vspace{.4cm}
(b) Show that $K(A,B)=P(A \cap B)- P(A)P(B)$ is also pos. definite.\\
\vspace{.4cm}
Let $A=(A_1, A_2, ... A_n)$ be a vector of events.
$\sum_{ij} \alpha_i \alpha_j\left( P(A \cap B)- P(A)P(B)\right) =\sum_{ij} \alpha_i \alpha_j E(\mathbb{I}_{A \cup  B})=$\\
$\sum_{ij} \alpha_i \alpha_j \left(E(\mathbb{I}_{A_i} \mathbb{I}_{A_j} )- E(\mathbb{I}_{A_i}) E(\mathbb{I}_{A_j})\right)=$\\
$\sum_{ij} \alpha_i \alpha_j Cov \left(\mathbb{I}_{A_i}(\omega),\mathbb{I}_{A_j}(\omega)\right)$\\
\vspace{.4cm}
where $\mathbb{I}_{A}(\omega)=1$ if $\omega\in A$ and $\mathbb{I}_{A}(\omega)= 0$, otherwise. Furthermore\\
\vspace{.4cm}
$\sum_{ij} \alpha_i \alpha_j Cov \left(\mathbb{I}_{A_i}(\omega),\mathbb{I}_{A_j}(\omega)\right)=$\\
$\sum \alpha_i^2 Var(\mathbb{I}_{A_i}(\omega))+2 \sum_{i,j|i<j} \alpha_i \alpha_j Cov\left(\mathbb{I}_{A_i}(\omega),\mathbb{I}_{A_j}(\omega)\right)=$\\
$\sum_i^n Var(\mathbb{I}_{A_i}(\omega)) \geq 0$.\\
\vspace{.4cm}
3. $\mathcal{X} =\mathbb{N}^{+}$\\
\vspace{.4cm}
(a) Show that $K(x,y)=LCM(x,y)$ is not pos. definite on $\mathbb{N}^{+}$.\\
\vspace{.4cm}
Pick $\mathcal{X} = \{ x_1,x_2\}$ where the $x_i$ are distinct points in $\mathbb{N}^{+}$.\\
Choose $(\alpha_1,\alpha2) \in \mathbb{R}^2$ with $\alpha_1=-\alpha2$. \\
\vspace{.4cm}
Write $\{x_p\}_k=\{n x_k\}$ for $n \in \mathbb{N}$. Then \\
\vspace{.4cm}
$\sum_{ij}^n \alpha_1 \alpha_j K(x_i,x_j)=\sum_{ij}^n \alpha_1 \alpha_j \min(\{x_p\}_i \cap \{x_p\}_j)=$\\
$=\alpha_1^2 \min\{x_p\}_1-2 \alpha_1^2 \min(\{x_p\}_i \cap \{x_p\}_j)+\alpha_1^2 \min\{x_p\}_2=$\\
$=\alpha_1^2\left(\min\{x_p\}_1+\min\{x_p\}_2-2 \min(\{x_p\}_i \cap \{x_p\}_j)\right) \leq 0$\\
\vspace{.4cm}
(b) Show that $K(x,y)=GCD(x,y)$ is pos. definite on $\mathbb{N}^{+}$.\\
\vspace{.5cm}
Proof by Beslin and Ligh (1989):
Let $X=(x_1,...x_n) \in \mathbb{N}^{n+}, D=\{d_1,...d_m\}$\\
be the set X union all the divisors of $x_i \in X$.\\
\vspace{.5cm}
Let $\phi(x)$ denote the Euler's totient function. If we define a new $n \times m$ matrix\\
$A=(a_{ij})$ and set\\
\vspace{.4cm}
\hspace{4cm} $a_{ij}:=e_{ij} \cdot \sqrt{\phi{d_j}}$ , with
\vspace{.4cm}

   \hspace{4cm}         $e_{ij}=$
                    $\left\{1 \hspace{.2cm} \text{if} \hspace{.2cm} d_j| x_i,\hspace{.2cm} 0 \hspace{.2cm} \text{otherwise} \right.$

\vspace{.4cm}
Then we have\\
\vspace{.4cm}
$\left(A A^{\top}\right)_{ij} = \sum_k^m a_{ik} a_{jk}=\sum \limits_{d_k|x_i, d_k|x_j} \sqrt{\phi(d_k)}\sqrt{\phi(d_k)}=\sum\limits_{d_k|gcd(x_i,x_j)} \phi(d_k)=$\\
$GCD(x_i,x_j)$\\
\vspace{.4cm}
The last equality is a known property of the totient function.\\
\vspace{.4cm}
Since any matrix that can be decocmposed into the product of another matrix\\
with its transpose is pos. definite, we can conclude that $K(x,y)$ is as well.\\
\vspace{.4cm}
4. let $x,y\in \mathbb{R}$ and let $k_\sigma(x,y)=\exp \left(-\frac{(x-y)^2}{2 \sigma^2}\right)$\\
\vspace{.4cm}
(a) For $\sigma=1$. Prove that $K(x,y)$ is pos.definite using the the characteristic function.\\
\vspace{.4cm}
Since E[$e^{itZ}$] = $e^{-t^2/2}$ we can set $x_i-x_j=t$ and notice the characteristic function\\
is symmetric with respect to $x_i,x_j$. Then for $(\alpha_1...\alpha_n)\in\mathbb{R}^n$\\
\vspace{.4cm}
$\sum_{ij}^{n} k(x_i,x_j)=\sum_{ij}^{n} E \left(e^{(x_i-x_j)Z}\right) = E\left(\sum_{ij}^{n}(e^{iZx_i} e^{-iZx_j})\right)=$\\
$E\left(\sum_{i}^{n}e^{iZx_i}\sum_{j}^{n} e^{-iZx_j}\right)=E\left(\sum_k \alpha_k e^{iZx_k}\right)^2 \geq 0$\\
\vspace{.4cm}
The last equality follows from the symmetry of K.\\
(b) Show that $k_\sigma$ is pos. definite for any $\sigma >0$\\
\vspace{.4cm}
Set $t=(x_i-x_j)/\sigma$, then following part (a) exactly\\
$E\left(\sum_{i}^{n}e^{iZx_i/\sigma}\sum_{j}^{n} e^{-iZx_j/\sigma}\right)=E\left(\sum_k \alpha_k e^{iZx_k/\sigma}\right)^2 \geq 0$\\



\end{flushleft}
\end{document}
